
The existence of a Quine is surprising because it would appear that a program that prints its own source would need to have the following 3 distinct portions of code:

\begin{itemize}
\item $S$ for some representation of the source code of the Quine.
\item $R$ for the code needed to read $S$.
\item $P$ for the code needed to print $S$.
\end{itemize}
So that somehow $S = R + S + P$ is satisfied, which seems paradoxical.

The existence of Quines (in any general purpose programming language - or more generally, any Turing complete model of computation) follows form Kleene's Recursion Theorem.

Specifically, let $Q: \mathbb{N} \times \mathbb{N} \rightarrow \mathbb{N}$ be the computable function defined by mapping $(x, y) \mapsto x$. By Kleene's Recursion Theorem, there exists a program $q$ such that:
\begin{equation*}
    \forall\ y \in \mathbb{N}:\ q(y) = Q(source(q), y) = source(q)
\end{equation*}
Hence, $q$ is a Quine, and importantly, a product of a fixed-point theorem: Kleene's Recursion Theorem.